% ============================================================================
%  MailForge 课程设计汇报 Beamer 幻灯片（朴素但不单调版）
%  编译：xelatex MailForge_presentation_20p.tex 连续两次
% ============================================================================
\documentclass[aspectratio=169,10pt]{beamer}
\usepackage[UTF8]{ctex}
\usepackage{fontspec}
\usepackage{amsmath,amssymb}
\usepackage{booktabs}
\usepackage{array}
\usepackage{multirow}
\usepackage{colortbl}
\usepackage{listings}
\usepackage{tikz}
\usetikzlibrary{arrows.meta,positioning,shapes.geometric,fit,calc,backgrounds}
\usepackage{graphicx}
\usepackage{xcolor}
\usepackage{hyperref}
\usetheme{default}
\usecolortheme{default}
\setbeamertemplate{navigation symbols}{}
\setbeamertemplate{caption}[numbered]
\setbeamertemplate{itemize items}[circle]
\setbeamertemplate{enumerate items}[default]
\setbeamertemplate{blocks}[rounded][shadow=false]

% ---------------- 克制的低饱和配色 ----------------
\definecolor{Paper}{RGB}{250,249,246}
\definecolor{Ink}{RGB}{38,46,55}
\definecolor{Accent}{RGB}{31,111,139}
\definecolor{Sage}{RGB}{92,141,137}
\definecolor{Terracotta}{RGB}{188,105,74}
\definecolor{Ochre}{RGB}{198,150,72}
\definecolor{PaperLine}{RGB}{226,230,233}
\definecolor{MFBlue}{RGB}{31,111,139}
\definecolor{MFAccent}{RGB}{45,128,154}
\definecolor{MFGreen}{RGB}{92,141,137}
\definecolor{MFRed}{RGB}{188,105,74}
\definecolor{MFOrange}{RGB}{198,150,72}
\definecolor{MFGray}{RGB}{247,246,242}
\definecolor{MFDark}{RGB}{38,46,55}

\setbeamercolor{background canvas}{bg=Paper}
\setbeamercolor{normal text}{fg=Ink}
\setbeamercolor{structure}{fg=Accent}
\setbeamercolor{title}{fg=Ink}
\setbeamercolor{subtitle}{fg=Accent}
\setbeamercolor{frametitle}{fg=Ink,bg=Paper}
\setbeamercolor{block title}{fg=white,bg=Accent}
\setbeamercolor{block body}{bg=PaperLine!35}
\setbeamercolor{alerted text}{fg=Terracotta}
\setbeamercolor{example text}{fg=Sage}
\setbeamercolor{footline}{fg=Ink!65,bg=Paper}
\setbeamerfont{title}{size=\LARGE,series=\bfseries}
\setbeamerfont{subtitle}{size=\large}
\setbeamerfont{frametitle}{size=\large,series=\bfseries}
\setbeamerfont{block title}{size=\small,series=\bfseries}
\setbeamerfont{footline}{size=\tiny}

% ---------------- 标题样式：左侧细色条 ----------------
\setbeamertemplate{frametitle}{%
  \vspace{0.12em}
  \begin{beamercolorbox}[wd=\paperwidth,ht=0.72cm,dp=0.06cm,leftskip=0.78cm]{frametitle}
    {\color{Accent}\rule{0.08cm}{0.50cm}}\hspace{0.32cm}%
    \usebeamerfont{frametitle}\insertframetitle
  \end{beamercolorbox}
  \vspace{-0.31em}
  {\color{Accent!30}\hspace{0.78cm}\rule{\dimexpr\paperwidth-1.56cm\relax}{0.6pt}}
}

% ---------------- 页脚：细线 + 页码 ----------------
\setbeamertemplate{footline}{%
  \vspace{0.08em}%
  \hbox to \textwidth{\usebeamerfont{footline}\insertshorttitle\hfill\insertframenumber\,/\,\inserttotalframenumber}%
  \vspace{-0.12em}%
  {\color{Accent!22}\rule{\textwidth}{0.5pt}}%
}

% ---------------- 标题页：留白 + 淡色块 ----------------
\setbeamertemplate{title page}{%
  \begin{tikzpicture}[remember picture,overlay]
    \fill[Accent!7] (current page.south west) rectangle ([yshift=6.0cm]current page.south east);
    \fill[Sage!8] ([xshift=9.2cm,yshift=0.9cm]current page.south west) rectangle ([xshift=10.6cm,yshift=5.8cm]current page.south west);
    \node[anchor=west,font=\Huge\bfseries,text=Ink]
      at ([xshift=1.1cm,yshift=2.15cm]current page.south west) {\inserttitle};
    \node[anchor=west,font=\large,text=Accent]
      at ([xshift=1.1cm,yshift=1.35cm]current page.south west) {\insertsubtitle};
    \draw[Accent!45,line width=0.8pt]
      ([xshift=1.1cm,yshift=0.95cm]current page.south west) -- ([xshift=8.2cm,yshift=0.95cm]current page.south west);
  \end{tikzpicture}
  \vspace{6.8cm}
  \begin{center}
    {\small\color{Ink!70}\insertauthor\quad\textbar\quad\insertinstitute}\\[0.25em]
    {\small\color{Ink!50}\insertdate}
  \end{center}
}

% ---------------- 代码与命令样式 ----------------
\newcommand{\key}[1]{\textcolor{Accent}{\textbf{#1}}}
\newcommand{\warn}[1]{\textcolor{Terracotta}{\textbf{#1}}}
\newcommand{\good}[1]{\textcolor{Sage}{\textbf{#1}}}
\newcommand{\code}[1]{\texttt{\small\detokenize{#1}}}

\lstdefinestyle{cppstyle}{
  language=C++,
  basicstyle=\ttfamily\scriptsize,
  keywordstyle=\color{Accent}\bfseries,
  commentstyle=\color{Sage!90!black},
  stringstyle=\color{Terracotta!85!black},
  numbers=left,
  numberstyle=\tiny\color{Ink!35},
  numbersep=5pt,
  frame=single,
  rulecolor=\color{PaperLine},
  backgroundcolor=\color{PaperLine!25},
  breaklines=true,
  columns=fullflexible,
  showstringspaces=false,
  tabsize=2,
  xleftmargin=1.2em
}
\lstdefinestyle{protostyle}{
  basicstyle=\ttfamily\scriptsize,
  frame=single,
  rulecolor=\color{PaperLine},
  backgroundcolor=\color{PaperLine!25},
  breaklines=true,
  columns=fullflexible,
  showstringspaces=false,
  xleftmargin=1.2em
}

\title[MailForge]{MailForge 邮件服务器}
\subtitle{SMTP + POP3 + Web + 两层加密 \enspace|\enspace 课程设计汇报}
\author[苏俊玮、王诗贺]{苏俊玮、王诗贺}
\institute[计算机与网络课程设计]{计算机与网络课程设计 · 选题 18}
\date{\today}

\begin{document}

% ============================== 1 封面 ==============================
\begin{frame}
  \titlepage
  \note{开场用一句话：MailForge 是从零实现的课程级邮件系统，协议、Web 和密码算法都尽量自己实现。}
\end{frame}

% ============================== 2 项目与需求 ==============================
\begin{frame}{项目简介与课程需求映射}
  \begin{columns}[T]
    \begin{column}{0.52\textwidth}
      \begin{block}{MailForge 是什么}
        \small
        \begin{itemize}
          \item 原生 Socket 手写 SMTP / POP3 / HTTP。
          \item 多用户 Web 邮箱：注册、登录、收发、附件、已发送。
          \item 邮件按 \code{.eml} 落盘，重启不丢。
          \item 两层加密：Web RSA 信封 + 服务器内部/落盘自研 AES/ChaCha。
          \item 内建协议演示终端与 100 封压测。
        \end{itemize}
      \end{block}
      \begin{block}{运行约定}
        \small
        SMTP: 2525 \quad POP3: 1110 \quad HTTP: 8080 \\
        默认账号：\code{alice/123456}，\code{bob/123456}
      \end{block}
    \end{column}
    \begin{column}{0.46\textwidth}
      \begin{block}{课程要求对照}
        \small
        \begin{tabular}{p{0.42\textwidth}p{0.48\textwidth}}
          \toprule
          要求 & 实现 \\
          \midrule
          SMTP 标准交互 & EHLO/MAIL/RCPT/DATA/QUIT \\
          POP3 标准交互 & USER/PASS/STAT/LIST/RETR/DELE/QUIT \\
          正文与附件 & MIME multipart + Base64 \\
          传输加密 & 两层加密，算法可切换 \\
          $\ge$ 2 种算法 & AES-256-CBC、ChaCha20 \\
          时延/成功率 & 100 封压测，内建统计 \\
          \bottomrule
        \end{tabular}
      \end{block}
      \begin{alertblock}{汇报口径}
        \small
        协议正确性用 RFC 状态机和端到端测试证明；密码算法用 NIST/RFC 官方向量回归证明。
      \end{alertblock}
    \end{column}
  \end{columns}
  \note{这一页交代项目全貌和验收对照，不要展开代码。}
\end{frame}
% ============================== 3 总体架构 ==============================
\begin{frame}{总体架构：三服务器 + 两客户端 + 两加密层}
  \centering
  \begin{tikzpicture}[
      node distance=0.7cm and 0.9cm,
      box/.style={draw, rounded corners, thick, align=center, minimum height=0.7cm, minimum width=2.0cm, font=\scriptsize},
      svr/.style={box, fill=MFBlue!8},
      cli/.style={box, fill=MFGreen!10},
      store/.style={box, fill=MFOrange!12},
      layer1/.style={draw=MFRed, dashed, rounded corners, inner sep=5pt},
      layer2/.style={draw=MFBlue, dashed, rounded corners, inner sep=5pt},
      arrow/.style={-{Latex[length=2mm]}, thick}
    ]
    \node[box, fill=MFGray] (browser) {浏览器\\WebCrypto};
    \node[svr, right=of browser] (http) {HttpServer\\:8080};
    \node[cli, right=of http] (clients) {SmtpClient\\Pop3Client};
    \node[svr, right=of clients] (mail) {SMTP :2525\\POP3 :1110};
    \node[store, below=0.6cm of mail] (mailbox) {mailbox/<user>/*.eml\\sent/<user>/*.eml};
    \node[store, below=0.6cm of clients] (keys) {keys/\\<user>.key\\server\_*.pem};
    \draw[arrow] (browser) -- node[above,font=\tiny]{HTTP/REST} (http);
    \draw[arrow] (http) -- (clients);
    \draw[arrow] (clients) -- node[above,font=\tiny]{SMTP/POP3} (mail);
    \draw[arrow] (mail) -- (mailbox);
    \draw[arrow] (http) -- (keys);
    \draw[arrow] (mail) -- (keys);
    \begin{scope}[on background layer]
      \node[layer1, fit=(browser)(http)] {};
      \node[layer2, fit=(http)(mail)(mailbox)(keys)] {};
    \end{scope}
    \node[MFRed, font=\tiny] at ($(browser)!0.5!(http)+(0,1.05)$) {层1：RSA 信封};
    \node[MFBlue, font=\tiny] at ($(http)!0.5!(mail)+(0,1.05)$) {层2：自研对称加密};
  \end{tikzpicture}
  \vspace{0.3em}
  \begin{itemize}
    \small
    \item 三个服务器各占一个线程 accept；每个连接再派发一个工作线程。
    \item HTTP 后端不直接写文件，而是作为 SMTP/POP3 客户端走本机协议，保证与外部协议客户端一致。
  \end{itemize}
  \note{重点讲两个虚线框：层1保护浏览器到服务器，层2保护服务器内部与磁盘。}
\end{frame}

% ============================== 4 目录与模块 ==============================
\begin{frame}[fragile]{目录结构与模块职责}
  \begin{columns}[T]
    \begin{column}{0.48\textwidth}
      \begin{lstlisting}[style=protostyle,basicstyle=\ttfamily\tiny]
MailForge/
+-- MailServer/
|   +-- Server.*       网络基类
|   +-- SmtpServer.*   SMTP 服务端
|   +-- Pop3Server.*   POP3 服务端
|   +-- SmtpClient.*   内部发信
|   +-- Pop3Client.*   内部收信
|   +-- HttpServer.*   Web 后端
|   +-- MailCrypto.*   加密统一入口
|   +-- web/           前端页面
+-- crypto/
|   +-- aes.* chacha20.*
|   +-- rsa.* random.*
|   +-- tests/ docs/
+-- common/
    +-- base64 / logger / file_util / socket
      \end{lstlisting}
    \end{column}
    \begin{column}{0.48\textwidth}
      \begin{block}{分层设计}
        \small
        \begin{enumerate}
          \item 公共层：Socket、Base64、日志、文件工具。
          \item 协议层：SMTP/POP3 服务端与客户端状态机。
          \item Web 层：HTTP 解析、REST、会话、MIME、压测。
          \item 密码层：MailCrypto 入口；自研 AES/ChaCha + OpenSSL RSA。
          \item 展示层：Web 邮箱、协议演示终端。
        \end{enumerate}
      \end{block}
      \begin{block}{关键解耦}
        \small
        SMTP/POP3 服务器只处理原文，不感知加密；加密挂在 HTTP 发送前和读取后。
      \end{block}
    \end{column}
  \end{columns}
  \note{这一页回答代码在哪里。重点强调协议层与密码层解耦。}
\end{frame}
% ============================== 5 SMTP 原理 ==============================
\begin{frame}[fragile]{SMTP 原理与标准交互}
  \begin{columns}[T]
    \begin{column}{0.52\textwidth}
      \begin{itemize}
        \item SMTP 是基于 TCP 的文本命令/响应协议，服务器问候码 220。
        \item 核心流程：
          \begin{enumerate}
            \item \code{EHLO/HELO}；
            \item \code{MAIL FROM:<...>} 信封发件人；
            \item \code{RCPT TO:<...>} 信封收件人；
            \item \code{DATA} 开始传正文，单独一行点结束；
            \item \code{QUIT} 退出。
          \end{enumerate}
        \item 2xx 成功、3xx 中间态、4xx 暂时失败、5xx 永久失败。
        \item 本项目实现核心命令，未实现 AUTH、STARTTLS、多收件人队列。
      \end{itemize}
    \end{column}
    \begin{column}{0.46\textwidth}
      \begin{block}{SMTP 对话}
        \tiny
        \begin{lstlisting}[style=protostyle,basicstyle=\ttfamily\tiny]
S: 220 MyMailServer ESMTP ready
C: EHLO client
S: 250 Hello client
C: MAIL FROM:<alice@example.com>
S: 250 OK
C: RCPT TO:<bob@example.com>
S: 250 OK
C: DATA
S: 354 End data with <CR><LF>.<CR><LF>
C: Subject: Hello
C:
C: body
C: .
S: 250 Message accepted
C: QUIT
S: 221 Bye
        \end{lstlisting}
      \end{block}
    \end{column}
  \end{columns}
  \note{强调信封地址与邮件头部 From/To 是两个概念。}
\end{frame}

% ============================== 6 SMTP 实现与落盘 ==============================
\begin{frame}[fragile]{SMTP 状态机、点填充与落盘}
  \begin{columns}[T]
    \begin{column}{0.46\textwidth}
      \begin{tikzpicture}[
          node distance=0.5cm and 0.6cm,
          st/.style={draw, rounded corners, thick, fill=MFBlue!8, align=center, minimum height=0.6cm, font=\tiny},
          arrow/.style={-{Latex[length=1.5mm]}, thick}
        ]
        \node[st] (a) {220};
        \node[st, right=of a] (b) {EHLO};
        \node[st, right=of b] (c) {MAIL};
        \node[st, right=of c] (d) {RCPT};
        \node[st, below=of c] (e) {DATA/354};
        \node[st, right=of e] (f) {单独点\\保存};
        \draw[arrow] (a)--(b); \draw[arrow] (b)--(c); \draw[arrow] (c)--(d);
        \draw[arrow] (d)|-(e); \draw[arrow] (e)--(f);
      \end{tikzpicture}
      \begin{itemize}
        \small
        \item \code{dataMode} 区分命令与正文。
        \item \code{SmtpMail} 跨命令保存会话。
        \item 点填充：正文行以点开头写成两个点，服务器还原。
      \end{itemize}
    \end{column}
    \begin{column}{0.50\textwidth}
      \begin{block}{落盘与多用户隔离}
        \small
        \begin{enumerate}
          \item 从收件人取 \code{@} 前部分，转小写；
          \item 存入 \code{mailbox/<user>/时间戳_随机数.eml}；
          \item 补齐缺失的 Date/From/To/Subject；
          \item 头部区 + 空行 + 正文写入文件。
        \end{enumerate}
      \end{block}
      
    \end{column}
  \end{columns}
  \note{落盘采用标准 .eml，协议客户端和 Web 后端读写同一份数据。}
\end{frame}
% ============================== 7 POP3 原理 ==============================
\begin{frame}[fragile]{POP3 原理：三阶段状态机与延迟删除}
  \begin{columns}[T]
    \begin{column}{0.52\textwidth}
      \begin{itemize}
        \item 服务器先回 \code{+OK}，出错回 \code{-ERR}。
        \item 三阶段：
          \begin{enumerate}
            \item AUTHORIZATION：\code{USER / PASS}；
            \item TRANSACTION：\code{STAT / LIST / RETR / DELE / RSET / NOOP}；
            \item UPDATE：\code{QUIT} 时真正删除被标记邮件。
          \end{enumerate}
        \item \code{DELE} 只打标记，中途断开不会误删。
        \item 多行响应用单独一行点结束，正文点开头行同样点填充。
        \item 支持新注册账号即时登录：内存表没有则重读 \code{users.txt}。
      \end{itemize}
    \end{column}
    \begin{column}{0.46\textwidth}
      \begin{block}{POP3 对话}
        \tiny
        \begin{lstlisting}[style=protostyle,basicstyle=\ttfamily\tiny]
S: +OK MailForge POP3 ready
C: USER bob
S: +OK User bob known
C: PASS 123456
S: +OK mailbox ready, 2 message(s)
C: STAT
S: +OK 2 12345
C: LIST
S: +OK 2 message(s)
S: 1 6000
S: 2 6345
S: .
C: RETR 1
S: +OK 6000 octets
... mail ...
S: .
C: DELE 1
S: +OK message 1 deleted
C: QUIT
S: +OK signing off
        \end{lstlisting}
      \end{block}
    \end{column}
  \end{columns}
  \note{重点讲 DELE 加 QUIT 的延迟删除语义。}
\end{frame}

% ============================== 8 POP3 实现 ==============================
\begin{frame}{POP3 实现：快照 + 会话状态}
  \begin{columns}[T]
    \begin{column}{0.50\textwidth}
      \begin{block}{关键结构}
        \small
        \begin{itemize}
          \item \code{Pop3State}：认证标志、用户名、\code{userKey}、邮件快照。
          \item \code{Pop3Mail}：路径、大小、\code{deleted} 标记。
          \item 登录成功后扫描 \code{mailbox/<user>/*.eml}，排序生成稳定编号。
        \end{itemize}
      \end{block}
      \begin{block}{认证与热加载}
        \small
        \code{normalizeUser()} 统一小写并去掉域名；\code{isUserKnown/checkPassword} 自动重读账号文件；登录成功同时 \code{ensureUserKey()}。
      \end{block}
    \end{column}
    \begin{column}{0.48\textwidth}
      \begin{itemize}
        \item \code{sendMailContent()} 逐行读文件，处理换行与点填充，不一次性载入大邮件。
        \item \code{QUIT} 时才 \code{unlink()} 被标记文件。
        \item HTTP 后端调用 \code{Pop3Client}，与外部邮件客户端同一套协议。
      \end{itemize}
      
    \end{column}
  \end{columns}
  \note{快照加标记：会话开始时读列表，DELE 只改内存标记，QUIT 扫描删除。}
\end{frame}
% ============================== 9 Server 与 HTTP ==============================
\begin{frame}[fragile]{网络基类与 Web 后端}
  \begin{columns}[T]
    \begin{column}{0.48\textwidth}
      \begin{block}{Server 基类}
        \small
        \code{socket -> bind -> listen -> accept}，每个连接一个线程，子类实现 \code{handleClient}。SMTP/POP3/HTTP 共用这一套骨架。
      \end{block}
      \begin{lstlisting}[style=cppstyle,basicstyle=\ttfamily\tiny]
class Server {
protected:
  virtual void handleClient(int fd) = 0;
};
      \end{lstlisting}
      
    \end{column}
    \begin{column}{0.50\textwidth}
      \begin{block}{HTTP 后端}
        \small
        \begin{itemize}
          \item 逐行解析请求行、查询串、头部、\code{Content-Length} body。
          \item POST 表单按 \code{a=b\&c=d} 解析并 URL 解码。
          \item 响应固定 \code{Connection: close}。
          \item \code{/api/login} 用 POP3 试登录，成功生成 token 存入会话表。
        \end{itemize}
      \end{block}
      \begin{lstlisting}[style=cppstyle,basicstyle=\ttfamily\tiny]
void HttpServer::handleClient(int fd) {
  HttpRequest req; HttpResponse resp;
  if (!readRequest(fd, req)) return;
  route(req, resp);
  sendHttp(fd, resp);
}
      \end{lstlisting}
    \end{column}
  \end{columns}
  \note{HTTP 是最小可用实现；强调它只是浏览器的翻译官，真正收发仍走 SMTP/POP3。}
\end{frame}

% ============================== 10 REST API 与 MIME ==============================
\begin{frame}{REST API、MIME 附件与统一 .eml 存储}
  \scriptsize
  \begin{columns}[T]
    \begin{column}{0.48\textwidth}
      \begin{block}{主要接口}
        \begin{tabular}{lp{0.52\textwidth}}
          \toprule
          路径 & 功能 \\
          \midrule
          \code{/api/register} & 注册 \\
          \code{/api/login} & POP3 验证 + token \\
          \code{/api/send} & 发信；可选 AES/ChaCha；可选 RSA 信封 \\
          \code{/api/inbox} & POP3 LIST + RETR \\
          \code{/api/mail?n=} & 读信、自动解密、解析附件 \\
          \code{/api/attachment} & 下载附件；\code{web=1} 走信封 \\
          \code{/api/delete} & DELE + QUIT 删除 \\
          \code{/api/sent} & 已发送列表 / 阅读 / 附件 \\
          \code{/api/benchmark} & 100 封压测：plain/aes/chacha/all \\
          \code{/api/webkey} & 下发服务器 RSA 公钥 \\
          \code{/api/webpub} & 上传浏览器 RSA 公钥 \\
          \bottomrule
        \end{tabular}
      \end{block}
    \end{column}
    \begin{column}{0.50\textwidth}
      \begin{block}{MIME 附件}
        \small
        \begin{itemize}
          \item 生成随机 boundary，组装 \code{multipart/mixed}。
          \item 附件段使用 Base64。
          \item 加密时整个 multipart 先整体加密，再作为正文发送。
          \item 下载时先解密邮件，再解析 MIME，最后 Base64 解码。
        \end{itemize}
      \end{block}
      \begin{block}{统一存储}
        \small
        SMTP 投递、POP3 读取、Web 收发围绕同一份 \code{.eml}；用户目录隔离，重启不丢。
      \end{block}
    \end{column}
  \end{columns}
  \note{API 表不逐条念，重点标发送、收件、附件、压测、Web RSA 交换。}
\end{frame}
% ============================== 11 两层加密 ==============================
\begin{frame}{两层加密模型：信任边界不同}
  \centering
  \begin{tikzpicture}[
      node distance=0.5cm and 1.0cm,
      box/.style={draw, rounded corners, thick, align=center, minimum height=0.8cm, minimum width=2.1cm, font=\small},
      arrow/.style={-{Latex[length=2.2mm]}, thick}
    ]
    \node[box, fill=MFGray] (browser) {浏览器\\WebCrypto};
    \node[box, fill=MFRed!12, right=1.2cm of browser] (http) {HTTP 服务器\\OpenSSL};
    \node[box, fill=MFGreen!12, right=1.2cm of http] (proto) {SMTP/POP3\\服务器};
    \node[box, fill=MFOrange!15, right=1.2cm of proto] (disk) {磁盘 .eml\\密文};
    \draw[arrow, MFRed, line width=1pt] (browser) -- node[above,font=\tiny]{层1：AES-GCM + RSA-OAEP} (http);
    \draw[arrow, MFBlue, line width=1pt] (http) -- node[above,font=\tiny]{层2：AES-CBC / ChaCha20} (proto);
    \draw[arrow, MFBlue, line width=1pt] (proto) -- (disk);
  \end{tikzpicture}
  \vspace{0.6em}
  \begin{columns}[T]
    \begin{column}{0.48\textwidth}
      \begin{block}{层1：浏览器到服务器}
        \small
        \begin{itemize}
          \item 本机无 TLS 时保护 Web 表单、收件箱、附件。
          \item AES-256-GCM 加密内容，RSA-OAEP-2048 包装会话密钥。
          \item 浏览器 WebCrypto，服务器 OpenSSL EVP。
        \end{itemize}
      \end{block}
    \end{column}
    \begin{column}{0.48\textwidth}
      \begin{block}{层2：服务器内部与磁盘}
        \small
        \begin{itemize}
          \item 邮件正文和附件整体加密后再走 SMTP/POP3/落盘。
          \item AES-256-CBC 或 ChaCha20 可切换。
          \item 纯 C++ 自研，不依赖 OpenSSL。
          \item 发送副本用收件人密钥，已发送副本用发件人密钥。
        \end{itemize}
      \end{block}
    \end{column}
  \end{columns}
  \note{这是创新点核心。B/S 架构下服务器最终能解密，因此是分层机密性，不是严格端到端。}
\end{frame}

% ============================== 12 层1 RSA ==============================
\begin{frame}[fragile]{层1：RSA 数字信封与 WebCrypto 互通}
  \begin{columns}[T]
    \begin{column}{0.50\textwidth}
      \begin{block}{发送：浏览器到服务器}
        \small
        \begin{enumerate}
          \item 浏览器生成一次性 AES-256-GCM 会话密钥；
          \item 用服务器 RSA 公钥做 RSA-OAEP(SHA-256) 封装；
          \item 表单用 AES-GCM 加密，末尾带 16B tag；
          \item 服务器用私钥拆封，再解 GCM。
        \end{enumerate}
      \end{block}
      \begin{block}{反向：服务器到浏览器}
        \small
        浏览器生成 RSA-2048 密钥对，公钥上传 \code{/api/webpub}；服务器读信、附件、压测结果用该公钥封装返回。
      \end{block}
    \end{column}
    \begin{column}{0.48\textwidth}
      \begin{lstlisting}[style=protostyle,basicstyle=\ttfamily\tiny]
k=<Base64 RSA-OAEP-SHA256(32B session key)>
|iv=<Base64 12B GCM IV>
|ct=<Base64 AES-GCM(ciphertext || 16B tag)>
      \end{lstlisting}
      \begin{itemize}
        \small
        \item WebCrypto 导入 SPKI 公钥；JWK 私钥存 localStorage。
        \item 服务器 OpenSSL 读写 PEM，双方标准算法互通。
        \item WebCrypto 不可用时前端降级明文 HTTP。
      \end{itemize}
      
    \end{column}
  \end{columns}
  \note{讲混合加密：RSA 只包装会话密钥，AES-GCM 加密内容并提供完整性。}
\end{frame}
% ============================== 13 层2 算法 ==============================
\begin{frame}{层2：自研 AES-256-CBC 与 ChaCha20}
  \begin{columns}[T]
    \begin{column}{0.48\textwidth}
      \begin{block}{AES-256-CBC}
        \small
        \begin{itemize}
          \item 分组密码：块 16B，密钥 32B，14 轮。
          \item CBC 链接，PKCS\#7 填充并校验。
          \item 每封邮件随机 16B IV。
          \item 线格式：\code{MailForge::ENC::AES:: + Base64(IV || 密文)}。
        \end{itemize}
      \end{block}
    \end{column}
    \begin{column}{0.48\textwidth}
      \begin{block}{ChaCha20}
        \small
        \begin{itemize}
          \item 流密码：密钥 32B，nonce 12B，64B 密钥流块。
          \item ARX：加法、异或、循环移位，无查表。
          \item 20 轮；每封邮件随机 nonce。
          \item 线格式：\code{MailForge::ENC::CHA:: + Base64(nonce || 密文)}。
        \end{itemize}
      \end{block}
    \end{column}
  \end{columns}
  \vspace{0.4em}
  \begin{alertblock}{工程约定与安全边界}
    \small
    真实 Subject + 空行 + 正文/附件 multipart 整体加密；头部只留占位 Subject 与 \code{X-MailForge-Crypto}。AES-CBC 和 ChaCha20 本身都不提供完整性认证，建议升级为 AES-GCM 或 ChaCha20-Poly1305。
  \end{alertblock}
  \note{对比分组密码和流密码：AES-CBC 需要填充和 IV，ChaCha20 是流式异或。}
\end{frame}

% ============================== 14 AES 原理与实现 ==============================
\begin{frame}[fragile]{AES-256 原理与实现要点}
  \begin{columns}[T]
    \begin{column}{0.50\textwidth}
      \begin{block}{四个变换}
        \small
        \begin{enumerate}
          \item SubBytes：S 盒非线性替换；
          \item ShiftRows：按行循环左移；
          \item MixColumns：GF(2$^8$) 上乘固定矩阵；
          \item AddRoundKey：与轮密钥异或。
        \end{enumerate}
        AES-256 密钥扩展为 15 个轮密钥；最后一轮不做 MixColumns。
      \end{block}
      \begin{block}{S 盒不是抄表}
        \small
        在 GF(2$^8$) 中求逆 $a^{-1}=a^{254}$，再做仿射变换；逆 S 盒由正向映射反推。
      \end{block}
    \end{column}
    \begin{column}{0.48\textwidth}
      \begin{lstlisting}[style=cppstyle,basicstyle=\ttfamily\tiny]
bool Aes::Encrypt(key, iv, plaintext, ciphertext) {
  size_t pad = 16 - (plaintext.size() % 16);
  // PKCS#7 填充
  ExpandKey(key.data(), rk.data());
  for each block:
    blk = plaintext_block ^ prev;
    EncryptBlock(rk, blk, out);
    prev = out;
}
      \end{lstlisting}
      \begin{block}{正确性验证}
        \small
        使用 NIST SP 800-38A F.2.3 官方向量回归，逐块比对 AES-256-CBC 输出。
      \end{block}
    \end{column}
  \end{columns}
  \note{强调自研有官方向量验证，不是自己加密自己解密。}
\end{frame}
% ============================== 15 ChaCha 原理与实现 ==============================
\begin{frame}[fragile]{ChaCha20 原理与实现要点}
  \begin{columns}[T]
    \begin{column}{0.52\textwidth}
      \begin{block}{状态与轮函数}
        \small
        16 个 32 位字：4 常量 + 8 密钥字 + 1 计数器 + 3 nonce 字。\\
        每轮做列轮与对角轮的 QuarterRound：
        \[
        a+=b;\ d\oplus=a;\ d\lll 16;\quad
        c+=d;\ b\oplus=c;\ b\lll 12
        \]
        共 10 次双轮，最后累加初始状态并输出 64B 密钥流。
      \end{block}
      \begin{block}{验证}
        \small
        使用 RFC 8439 §2.3.2 官方向量验证 \code{Block()}。
      \end{block}
    \end{column}
    \begin{column}{0.46\textwidth}
      \begin{lstlisting}[style=cppstyle,basicstyle=\ttfamily\tiny]
void ChaCha20::Block(key, nonce, counter, out) {
  uint32_t state[16], work[16];
  state[0..3] = constants;
  state[4..11] = key words;
  state[12] = counter;
  state[13..15] = nonce words;
  for (int round = 0; round < 10; ++round) {
    // column + diagonal QuarterRound
  }
  for (int i = 0; i < 16; ++i)
    Store32LE(work[i] + state[i], out + 4*i);
}
      \end{lstlisting}
      
    \end{column}
  \end{columns}
  \note{ChaCha20 的优势是无查表、软件实现友好；缺点是当前没有 Poly1305。}
\end{frame}

% ============================== 16 密钥与挂钩 ==============================
\begin{frame}{密钥管理与加解密挂钩点}
  \begin{columns}[T]
    \begin{column}{0.50\textwidth}
      \begin{block}{密钥文件}
        \small
        \begin{itemize}
          \item 每账号 \code{keys/<user>.key}：32B 随机密钥；
          \item 首次使用自动生成；
          \item 用户名规范化：取 \code{@} 前、转小写；
          \item 全局互斥锁保护查文件加写文件。
        \end{itemize}
      \end{block}
      \begin{block}{服务器 RSA 密钥}
        \small
        \code{keys/server_public.pem} / \code{server_private.pem}，首次启动自动生成。
      \end{block}
    \end{column}
    \begin{column}{0.48\textwidth}
      \begin{block}{加解密挂钩点}
        \small
        \begin{enumerate}
          \item \code{handleSend()}：用收件人密钥加密载荷；
          \item \code{decodeMail()}：按魔数识别算法，用查看者密钥解密；
          \item 已发送副本：用发件人密钥再加密一份；
          \item 附件下载：先解密邮件，再解析 MIME，再 Base64 解码。
        \end{enumerate}
      \end{block}
      
    \end{column}
  \end{columns}
  \note{零配置密钥是亮点，但安全级别是演示级。真实系统不能裸存对称密钥。}
\end{frame}
% ============================== 17 端到端与测试 ==============================
\begin{frame}{端到端数据流与测试结果}
  \centering
  \begin{tikzpicture}[
      node distance=0.35cm,
      flow/.style={draw, rounded corners, thick, fill=MFGray, align=center, minimum height=0.6cm, text width=2.1cm, font=\tiny},
      arrow/.style={-{Latex[length=1.8mm]}, thick}
    ]
    \node[flow, fill=MFRed!10] (s1) {浏览器附件表单};
    \node[flow, fill=MFRed!10, right=of s1] (s2) {层1 AES-GCM\\+ RSA-OAEP};
    \node[flow, fill=MFBlue!10, right=of s2] (s3) {HTTP 解信封};
    \node[flow, fill=MFBlue!10, right=of s3] (s4) {MIME multipart};
    \node[flow, fill=MFBlue!10, below=of s4] (s5) {层2 AES/ChaCha};
    \node[flow, fill=MFGreen!10, left=of s5] (s6) {SMTP / .eml 密文};
    \node[flow, fill=MFGreen!10, left=of s6] (s7) {POP3 RETR};
    \node[flow, fill=MFBlue!10, left=of s7] (s8) {层2 解密 / MIME};
    \node[flow, fill=MFRed!10, below=of s8] (s9) {层1 信封返回};
    \draw[arrow] (s1)--(s2); \draw[arrow] (s2)--(s3); \draw[arrow] (s3)--(s4);
    \draw[arrow] (s4)--(s5); \draw[arrow] (s5)--(s6); \draw[arrow] (s6)--(s7);
    \draw[arrow] (s7)--(s8); \draw[arrow] (s8)--(s9);
  \end{tikzpicture}
  \vspace{0.5em}
  \begin{columns}[T]
    \begin{column}{0.50\textwidth}
      \begin{block}{正确性测试}
        \small
        \begin{itemize}
          \item AES-256-CBC：NIST SP 800-38A F.2.3。
          \item ChaCha20：RFC 8439 §2.3.2。
          \item RSA 信封：Seal/Open 双向、篡改和错密钥拒绝。
          \item 协议端到端：SMTP 发信 -> 落盘 -> POP3 收信。
        \end{itemize}
      \end{block}
    \end{column}
    \begin{column}{0.48\textwidth}
      \begin{block}{性能压测}
        \small
        每模式 100 封约 1MB 邮件：发送 100/100，丢包率 0\%，加密邮件解密还原 100/100；AES 平均约 100ms，ChaCha20 约 60ms，远低于 2s。
      \end{block}
      \begin{alertblock}{注意}
        \small
        以上为本地环回演示数据；真实网络评测应补充公网、并发、P95/P99 等指标。
      \end{alertblock}
    \end{column}
  \end{columns}
  \note{重点讲官方向量和端到端闭环，性能数据说明本地环回即可。}
\end{frame}

% ============================== 18 创新点 ==============================
\begin{frame}{创新点总结}
  \begin{columns}[T]
    \begin{column}{0.50\textwidth}
      \begin{block}{密码学创新}
        \small
        \begin{enumerate}
          \item 两层加密模型：Web 链路与服务器内部/落盘分别保护。
          \item 纯 C++ 自研 AES-256-CBC 与 ChaCha20，过 NIST/RFC 官方向量。
          \item WebCrypto 与 OpenSSL 互通，标准线格式。
          \item 加密与协议解耦，SMTP/POP3 对密文透明。
        \end{enumerate}
      \end{block}
    \end{column}
    \begin{column}{0.48\textwidth}
      \begin{block}{系统工程创新}
        \small
        \begin{enumerate}
          \setcounter{enumi}{4}
          \item 协议、Web、存储统一 \code{.eml}，不绕开协议。
          \item 浏览器协议演示终端 + 100 封压测 + 自动清理。
          \item 零配置密钥：首次使用自动生成。
          \item 已发送副本用发件人密钥再加密，双方都能读。
        \end{enumerate}
      \end{block}
      \begin{alertblock}{一句话}
        \small
        既展示了协议状态机，又展示了密码算法内部结构，并且都有可运行、可验证的闭环。
      \end{alertblock}
    \end{column}
  \end{columns}
  \note{创新点不需要每个都展开，重点突出两层加密、自研算法、WebCrypto 互通。}
\end{frame}
% ============================== 20 总结与问答 ==============================
\begin{frame}{总结与答辩准备}
  \begin{columns}[T]
    \begin{column}{0.48\textwidth}
      \begin{block}{总结}
        \small
        \begin{itemize}
          \item 从零手写 SMTP、POP3、HTTP。
          \item 多用户 Web 邮箱、MIME 附件、\code{.eml} 持久化。
          \item 两层加密：RSA 信封 + 自研 AES/ChaCha。
          \item 官方向量、端到端、100 封压测验证。
        \end{itemize}
      \end{block}
      \begin{block}{核心收获}
        \small
        协议要处理粘包/半包、点填充；密码算法要用官方向量证明；加密要按信任边界分层；安全加固与工程化同样重要。
      \end{block}
    \end{column}
    \begin{column}{0.48\textwidth}
      \begin{block}{高频问答}
        \small
        \begin{itemize}
          \item 粘包/半包：缓冲区按换行拆行，半行留到下次。
          \item POP3 删除：DELE 只标记，QUIT 才真正删除。
          \item 为什么两层：浏览器链路和服务器内部/磁盘信任边界不同。
          \item 自研算法安全吗：教学验证用，生产应用审计过的 AEAD。
        \end{itemize}
      \end{block}
      \begin{block}{谢谢}
        \centering
        MailForge：让邮件协议和安全算法都看得见、跑得通、验得过。
      \end{block}
    \end{column}
  \end{columns}
  \note{最后一页用于总结和提问。留出演示源码或 PDF 的时间。}
\end{frame}


\end{document}
